% Provenance (never printed): see sources/outline_sources.md (Year 3:
% partition coefficients; batch and fractional melting; crystallisation
% models; reading a spider diagram quantitatively). Brief:
% sources/book5/BRIEFS.md, ch. 12. Curves from figdata/bachelor-3/
% melting_models.py. Exercise and problem values in the Book 5 scratch
% calculator calc/ch12.py. Aldren Island and its lavas are invented.
\chapter{Pemodelan Unsur Jejak}\label{ch:b3:trace-element-modelling}

Dua basal meletus di pulau yang sama, berjarak beberapa kilometer. Yang satu
mengandung lantanum lebih dari satu setengah kali lipat dibanding yang lain, dan hampir empat kali
lebih sedikit iterbium. Apakah keduanya meleleh dari dua mantel berbeda, atau dari satu mantel
yang meleleh dengan jumlah berbeda, atau pada kedalaman berbeda? Unsur mayor
hampir tidak memberi petunjuk: keduanya basal. Unsur jejak, yang mengikuti mineral
tanpa mengubahnya, menyimpan catatan tentang seberapa banyak lelehan dihasilkan, bagaimana
lelehan itu keluar dari sumbernya, dan mineral mana yang tertinggal. Bab ini menuliskan catatan itu
dalam bentuk persamaan dan membacanya kembali.

\begin{recall}
Koefisien partisi mineral--lelehan $D$, koefisien partisi total, unsur
kompatibel dan inkompatibel, unsur tanah jarang, kondrit CI
dan mantel primitif, diagram ternormalisasi, dan anomali europium,
MORB dan OIB ada di volume Tahun 2, dalam bab tentang unsur-unsur di
Bumi. Peridotit, lherzolit, pelelehan sebagian, dan fraksi lelehan $F$
ada di bab tentang pelelehan mantel. Kesetimbangan dan kristalisasi
fraksional, hukum fraksionasi Rayleigh, pencampuran dan
asimilasi--kristalisasi fraksional (AFC), dalam bab tentang diferensiasi
magmatik. Di sana telah dibuktikan bahwa rasio dua unsur yang sangat inkompatibel
tetap konstan selama kristalisasi.
\end{recall}

\section{Pembukuan batuan yang meleleh}\label{sec:b3:trace-element-modelling:balance}

Batuan bermassa 1 meleleh. Pada fraksi lelehan $F$, ia telah membentuk massa
$F$ lelehan dan menyisakan massa $1 - F$ residu padat. Suatu unsur jejak yang
hadir dengan konsentrasi $C_0$ dalam sumber terbagi di antara keduanya;
konsentrasinya dalam lelehan $C_L$, dalam residu $C_S$. Setiap model bab ini
adalah bentuk dari satu persamaan neraca, $C_0 = FC_L + (1 - F)C_S$
untuk seluruh lelehan dan residu, dan suatu aturan yang menghubungkan $C_S$
dengan $C_L$: pada kontak antara kristal dan lelehan, $C_{\text{mineral}} = D\,C_L$.
Model-model berbeda dalam hal apa yang bersentuhan dengan apa, dan kapan.

Koefisien partisi mineral-mineral mantel sangat bervariasi. Dalam satu
kompilasi, lantanum memiliki $D = 0.058$ dalam klinopiroksen dan $0.001$ dalam
garnet serta dalam olivin, ortopiroksen, dan spinel; iterbium memiliki $D = 0.58$ dalam
klinopiroksen, $4.5$ dalam garnet, dan $0.16$ pada ketiga lainnya. Garnet
menyerap unsur tanah jarang berat dan menolak yang ringan.
% ledger: budahn91:d-ree

\begin{definition}[Pelelehan non-modal]\label{def:b3:trace-element-modelling:nonmodal}
Pada \emph{pelelehan non-modal}\index{pelelehan non-modal}, mineral-mineral
sumber masuk ke dalam lelehan dengan proporsi $p_i$ yang berbeda dari fraksi
massanya $x_i$ dalam sumber. Lelehan kemudian melihat koefisien partisi total
sumber, $D_0 = \sum_i x_iD_i$, pada awalnya, dan koefisien berbobot pelelehan
$P = \sum_i p_iD_i$ seiring lelehan itu bertambah.
\end{definition}

Pelelehan jarang bersifat modal. Dalam satu contoh terhitung untuk lherzolit spinel dengan
\qty{58}{\percent} olivin, \qty{27}{\percent} ortopiroksen,
\qty{12}{\percent} klinopiroksen dan \qty{3}{\percent} spinel, mineral-mineral itu
masuk ke dalam lelehan dengan proporsi 20, 25, 45 dan \qty{10}{\percent}:
klinopiroksen meleleh empat kali lebih cepat daripada bagiannya, olivin tiga kali
lebih lambat.
% ledger: white13:ex73

\begin{proposition}[Koefisien partisi total residu]
\label{prop:b3:trace-element-modelling:residue}
Pada \omterm{def:b3:trace-element-modelling:nonmodal}{pelelehan non-modal} dengan $p_i$ konstan, residu pada fraksi lelehan $F$
memiliki koefisien partisi total
\[
\bar D(F) = \frac{D_0 - PF}{1 - F},
\]
berlaku selama setiap mineral masih ada, yaitu untuk $F < x_i/p_i$ untuk setiap $i$.
\end{proposition}

\begin{proof}
Per satuan massa sumber, massa mineral $i$ yang tersisa pada $F$ adalah $x_i - p_iF$,
dan residu bermassa $1 - F$. Koefisien totalnya adalah rata-rata berbobot massa,
$\sum_i (x_i - p_iF)D_i/(1 - F) = (D_0 - PF)/(1 - F)$. Suatu mineral habis
ketika $x_i - p_iF = 0$.
\end{proof}

Dengan koefisien mineral di atas, lherzolit spinel pada contoh terhitung dan
lherzolit garnet yang dipertahankan pada \qty{5}{GPa} dan
\qty{1400}{\celsius} di laboratorium, dengan \qty{54.1}{\percent} olivin,
\qty{9.8}{\percent} ortopiroksen, \qty{19.5}{\percent} klinopiroksen dan
\qty{16.7}{\percent} garnet menurut massa, diperoleh koefisien total berikut.
% ledger: white13:ex73, rfr25:gkr-5gpa, budahn91:d-ree
\begin{center}
\small
\begin{tabular}{lccc}
\hline
sumber & $D_{\mathrm{La}}$ & $D_{\mathrm{Sm}}$ & $D_{\mathrm{Yb}}$ \\
\hline
lherzolit spinel & 0.0078 & 0.062 & 0.21 \\
lherzolit garnet & 0.012 & 0.107 & 0.97 \\
\hline
\end{tabular}
\end{center}
Lantanum sangat inkompatibel pada keduanya. Iterbium cukup inkompatibel
pada lherzolit spinel dan hampir netral pada lherzolit garnet. Perbedaan
antara kedua sumber hampir seluruhnya terletak pada unsur tanah jarang
berat.

\section{Pelelehan tumpak}\label{sec:b3:trace-element-modelling:batch}

\begin{theorem}[Pelelehan tumpak]\label{thm:b3:trace-element-modelling:batch}
Pada \emph{pelelehan tumpak}\index{pelelehan tumpak}, lelehan tetap dalam kesetimbangan
dengan seluruh residu sampai lelehan keluar sekaligus dalam satu tumpak. Dengan
koefisien partisi total $D$ yang konstan, lelehan dan residu pada fraksi lelehan $F$
memiliki
\[
\frac{C_L}{C_0} = \frac{1}{D + F(1 - D)}, \qquad
\frac{C_S}{C_0} = \frac{D}{D + F(1 - D)}.
\]
\end{theorem}

\begin{proof}
Pada kesetimbangan, setiap kristal residu menyimpan $D_iC_L$, sehingga residu
rata-rata menyimpan $C_S = DC_L$. Kesetimbangan $C_0 = FC_L + (1 - F)DC_L$
memberi $C_L$, dan $C_S$ mengikuti.
\end{proof}

Rumusnya adalah hukum kristalisasi kesetimbangan dari volume Tahun 2 yang dibaca
dari belakang: pada keduanya, $F$ adalah fraksi cairan dalam sistem tertutup. Dua
batas perlu diketahui. Ketika $F \to 0$, $C_L \to C_0/D$: tetes lelehan
pertama paling banyak dapat memperkaya suatu unsur dengan faktor $1/D$. Ketika $D \to 0$, $C_L \to
C_0/F$: unsur yang benar-benar inkompatibel terencerkan ke dalam lelehan, dan
pengayaannya hanya bergantung pada $F$. Untuk unsur kompatibel, $D > 1$, lelehan
lebih miskin daripada sumbernya, dan makin miskin untuk $F$ yang makin kecil.

\begin{proposition}[Pelelehan tumpak non-modal]\label{prop:b3:trace-element-modelling:nonmodal}
Pada \omterm{thm:b3:trace-element-modelling:batch}{pelelehan tumpak} dengan proporsi pelelehan $p_i$ dari \cref{def:b3:trace-element-modelling:nonmodal},
\[
\frac{C_L}{C_0} = \frac{1}{D_0 + F(1 - P)}.
\]
\end{proposition}

\begin{proof}
Residu menyimpan $\bar D(F)\,C_L$ (\cref{prop:b3:trace-element-modelling:residue}),
sehingga $C_0 = FC_L + (1 - F)\bar D(F)C_L = C_L(F + D_0 - PF)$.
\end{proof}

Ketika $P$ melampaui $D_0$, seperti untuk lantanum dalam contoh terselesaikan, dengan $D_0 =
0.0074$ dan $P = 0.025$, residu kehilangan inang terbaiknya, klinopiroksen,
lebih cepat daripada mineral lainnya. Lelehan kemudian sedikit lebih kaya daripada
yang diramalkan pelelehan modal, sebesar kurang dari \qty{2}{\percent} pada $F = 0.1$ untuk lantanum (latihan~6), tetapi
residu kehabisan klinopiroksen pada $F = 0.12/0.45 \approx 0.27$, dan setelah
itu mineral dan koefisiennya berubah.
% ledger: white13:ex73

\section{Pelelehan fraksional dan dinamis}\label{sec:b3:trace-element-modelling:fractional}

Lelehan dalam peridotit berkumpul di sepanjang tepi butir dan dapat keluar begitu
saling terhubung. Anggota ujung yang berlawanan dengan \omterm{thm:b3:trace-element-modelling:batch}{pelelehan tumpak} menyingkirkan setiap
tetes begitu ia terbentuk.

\begin{theorem}[Pelelehan fraksional]\label{thm:b3:trace-element-modelling:fractional}
Pada \emph{pelelehan fraksional}\index{pelelehan fraksional}, setiap penambahan lelehan
berada dalam kesetimbangan dengan residu hanya pada saat ia terbentuk, lalu
meninggalkannya. Dengan $D$ konstan, setelah massa total $F$ lelehan terekstraksi,
residu dan lelehan yang terbentuk pada saat itu memiliki
\[
\frac{C_S}{C_0} = (1 - F)^{1/D - 1}, \qquad
\frac{C_L}{C_0} = \frac{1}{D}(1 - F)^{1/D - 1}.
\]
\end{theorem}

\begin{proof}
Residu memiliki massa $1 - F$ dan konsentrasi $C_S$. Ketika massa
$\mathrm dF$ lelehan keluar pada $C_L = C_S/D$, unsur dalam residu
berkurang sebesar $\mathrm d[(1 - F)C_S] = -(C_S/D)\,\mathrm dF$. Dengan menjabarkan,
$(1 - F)\,\mathrm dC_S = C_S(1 - 1/D)\,\mathrm dF$, sehingga
$\mathrm dC_S/C_S = (1/D - 1)\,\mathrm d[-\ln(1 - F)]$. Mengintegralkan dari
$C_S = C_0$ pada $F = 0$ memberi $\ln(C_S/C_0) = (1/D - 1)\ln(1 - F)$, dan
$C_L = C_S/D$.
\end{proof}

Untuk unsur inkompatibel, $1/D - 1$ besar dan positif: residu
terkuras dalam beberapa kali $D$ dari pelelehan, dan lelehan sesaat turun
dari $C_0/D$ hampir ke nol. Residu \omterm{thm:b3:trace-element-modelling:fractional}{pelelehan fraksional} jauh lebih miskin
daripada residu \omterm{thm:b3:trace-element-modelling:batch}{pelelehan tumpak}. Namun lelehan tidak dierupsikan satu per satu,
melainkan berkumpul dalam alur dan membentuk kolam; lelehan terkumpul, yang didefinisikan
di bawah, digambar putus-putus pada gambar.

\begin{omfigure}
\begin{tikzpicture}
\begin{groupplot}[group style={group size=2 by 1, horizontal sep=1.6cm},
  omrecord, width=0.47\linewidth, height=0.44\linewidth, title style={font=\small},
  xlabel={fraksi lelehan $F$}, xmin=0, xmax=0.3, ymode=log, ymin=0.01, ymax=150,
  xtick={0,0.1,0.2,0.3}, xticklabels={0,0.1,0.2,0.3},
  ytick={0.01,0.1,1,10,100}, yticklabels={0.01,0.1,1,10,100},
  legend style={font=\scriptsize, legend cell align=left}]
\nextgroupplot[title={(a) tumpak (padat), teragregasi (putus-putus)},
  ylabel={$C_L/C_0$}, legend style={at={(0.97,0.97)}, anchor=north east}]
% computed in figdata/bachelor-3/melting_models.py
\addplot[omDef, thick] table[x=f, y=b001] {figdata/out/bachelor-3-melting_models-curves.dat};
\addlegendentry{$D = 0.01$}
\addplot[red!70!black, thick] table[x=f, y=b01] {figdata/out/bachelor-3-melting_models-curves.dat};
\addlegendentry{$D = 0.1$}
\addplot[black, thick] table[x=f, y=b1] {figdata/out/bachelor-3-melting_models-curves.dat};
\addlegendentry{$D = 1$}
\addplot[omLithLime!70!black, thick] table[x=f, y=b10] {figdata/out/bachelor-3-melting_models-curves.dat};
\addlegendentry{$D = 10$}
\addplot[omDef, thick, densely dashed] table[x=f, y=a001] {figdata/out/bachelor-3-melting_models-curves.dat};
\addplot[red!70!black, thick, densely dashed] table[x=f, y=a01] {figdata/out/bachelor-3-melting_models-curves.dat};
\addplot[omLithLime!70!black, thick, densely dashed] table[x=f, y=a10] {figdata/out/bachelor-3-melting_models-curves.dat};
\nextgroupplot[title={(b) fraksional, sesaat}]
\addplot[omDef, thick] table[x=f, y=i001] {figdata/out/bachelor-3-melting_models-curves.dat};
\addplot[red!70!black, thick] table[x=f, y=i01] {figdata/out/bachelor-3-melting_models-curves.dat};
\addplot[black, thick] table[x=f, y=i1] {figdata/out/bachelor-3-melting_models-curves.dat};
\addplot[omLithLime!70!black, thick] table[x=f, y=i10] {figdata/out/bachelor-3-melting_models-curves.dat};
\end{groupplot}
\end{tikzpicture}

\par\smallskip{\small Pengayaan unsur jejak dalam lelehan terhadap
fraksi lelehan, untuk empat koefisien partisi total. (a) Lelehan tumpak (padat)
dan \omterm{def:b3:trace-element-modelling:fractional}{lelehan fraksional teragregasi} (putus-putus; untuk $D = 1$ keduanya tetap 1). (b) Lelehan
sesaat dari \omterm{thm:b3:trace-element-modelling:fractional}{pelelehan fraksional}: unsur dengan $D = 0.01$ habis
dari residu setelah beberapa persen pelelehan.}\label{fig:b3:trace-element-modelling:curves}
\end{omfigure}

\begin{definition}[Lelehan fraksional teragregasi]\label{def:b3:trace-element-modelling:fractional}
Yang disebut \emph{lelehan fraksional teragregasi}\index{lelehan fraksional teragregasi} adalah
campuran dari semua lelehan sesaat pada \omterm{thm:b3:trace-element-modelling:fractional}{pelelehan fraksional} dari $0$ hingga
$F$, yang dikumpulkan tanpa reaksi lebih lanjut dengan residu.
\end{definition}

\begin{proposition}[Komposisi lelehan teragregasi]
\label{prop:b3:trace-element-modelling:aggregate}
Dengan $D$ konstan, \omterm{def:b3:trace-element-modelling:fractional}{lelehan fraksional teragregasi} memiliki
\[
\frac{\bar C_L}{C_0} = \frac{1 - (1 - F)^{1/D}}{F}.
\]
\end{proposition}

\begin{proof}
Unsur yang diekstraksi adalah sumber dikurangi residu,
$C_0 - (1 - F)C_S = C_0[1 - (1 - F)^{1/D}]$ menurut \cref{thm:b3:trace-element-modelling:fractional},
tersebar pada massa $F$ lelehan.
\end{proof}

Begitu $(1 - F)^{1/D}$ kecil, lelehan teragregasi menampung seluruh
unsur dan $\bar C_L \approx C_0/F$, berapa pun $D$: untuk unsur yang sangat
inkompatibel, lelehan fraksional yang digabung dan \omterm{thm:b3:trace-element-modelling:batch}{pelelehan tumpak} dengan $D$ kecil sepakat, dan
keduanya mengukur $F$.

Pelelehan nyata berada di antara kedua ekstrem. Residu menyimpan sebagian lelehan di dalam
porinya, dan lelehan yang keluar adalah kelebihan di atas porositas itu.

\begin{definition}[Pelelehan dinamis]\label{def:b3:trace-element-modelling:dynamic}
Pada \emph{pelelehan dinamis}\index{pelelehan dinamis} (disebut juga pelelehan
kontinu), residu mempertahankan fraksi massa lelehan $\phi$ yang konstan di dalam
porinya, dalam kesetimbangan dengan kristalnya, dan setiap lelehan yang terbentuk melampaui itu
diekstraksi.
\end{definition}

\begin{proposition}[Pelelehan dinamis]\label{prop:b3:trace-element-modelling:dynamic}
\omterm{def:b3:trace-element-modelling:dynamic}{Pelelehan dinamis} dengan porositas $\phi$ adalah \omterm{thm:b3:trace-element-modelling:fractional}{pelelehan fraksional} dengan koefisien partisi efektif
$D' = (1 - \phi)D + \phi$, dengan $F$ adalah fraksi massa
lelehan yang diekstraksi:
\[
\frac{C_L}{C_0} = \frac{1}{D'}(1 - F)^{1/D' - 1}.
\]
\end{proposition}

\begin{proof}
Residu terdiri dari kristal dengan fraksi $1 - \phi$ dan lelehan pori dengan
fraksi $\phi$, keduanya dalam kesetimbangan dengan lelehan pada $C_L$, sehingga berlaku
$C_S = [(1 - \phi)D + \phi]C_L = D'C_L$. Pembuktian
\cref{thm:b3:trace-element-modelling:fractional} kemudian berlaku dengan $D'$.
\end{proof}

Untuk $\phi = 0$ model ini adalah \omterm{thm:b3:trace-element-modelling:fractional}{pelelehan fraksional}; untuk $\phi \to 1$ model ini
mendekati \omterm{thm:b3:trace-element-modelling:batch}{pelelehan tumpak}. Porositas satu persen membatasi pengayaan
lelehan pertama pada $1/\phi = 100$, berapa pun kecil $D$: lelehan pori berperan
sebagai zona penyangga bagi residu.

\begin{omfigure}
\begin{tikzpicture}[font=\small, >=Stealth]
\tikzset{rock/.style={draw, fill=omLithLime!25, minimum width=2.3cm, minimum height=2.0cm},
  melt/.style={draw, fill=red!25}}
% (a) batch
\node[rock] (a) at (0,0) {};
\foreach \x/\y in {-0.7/0.5, 0.3/0.6, -0.2/-0.1, 0.7/-0.4, -0.6/-0.6, 0.4/0.0}
  \fill[red!55] (\x,\y) circle (0.09);
\node[align=center, font=\scriptsize] at (0,-1.45) {residu $1-F$ pada $DC_L$\\lelehan $F$ pada $C_L$};
\draw[->, thick, red!70!black] (0,1.05) -- (0,1.7) node[above, font=\scriptsize] {sekaligus};
\node[font=\small] at (0,-2.25) {(a) tumpak};
% (b) fractional
\node[rock] (b) at (4,0) {};
\draw[->, thick, red!70!black] (4,1.05) -- (4,1.7);
\foreach \y/\s in {1.85/0.06, 2.05/0.08, 2.27/0.1}
  \fill[red!55] (4,\y) circle (\s);
\node[align=center, font=\scriptsize] at (4,-1.45) {setiap tetes keluar\\begitu terbentuk};
\node[font=\small] at (4,-2.25) {(b) fraksional};
% (c) dynamic
\node[rock] (c) at (8,0) {};
\foreach \x/\y in {7.4/0.5, 8.4/0.3, 7.8/-0.5, 8.6/-0.6}
  \fill[red!55] (\x,\y) circle (0.06);
\draw[->, thick, red!70!black] (8,1.05) -- (8,1.7);
\foreach \y/\s in {1.85/0.06, 2.05/0.08, 2.27/0.1}
  \fill[red!55] (8,\y) circle (\s);
\node[align=center, font=\scriptsize] at (8,-1.45) {pori menahan $\phi$;\\kelebihannya keluar};
\node[font=\small] at (8,-2.25) {(c) dinamis};
\end{tikzpicture}

\par\smallskip{\small Tiga cara lelehan keluar dari sumber yang meleleh. (a)
Tumpak: lelehan tetap bersama seluruh residu, lalu keluar. (b) Fraksional:
setiap penambahan keluar sekaligus; penambahan yang digabung membentuk
lelehan teragregasi. (c) Dinamis: residu menyimpan fraksi massa lelehan $\phi$ yang kecil di
dalam porinya.}\label{fig:b3:trace-element-modelling:modes}
\end{omfigure}

\section{Pelelehan, lalu kristalisasi}\label{sec:b3:trace-element-modelling:crystallisation}

Lelehan yang meninggalkan sumbernya jarang keluar ke permukaan tanpa berubah. Ia mengkristal
olivin, lalu piroksen dan plagioklas dalam perjalanan ke atas, ia mungkin mengasimilasi batuan
dinding, dan ia mungkin bercampur dengan kumpulan lain. Hukum-hukum volume Tahun 2 berlaku
secara bergantian: lava adalah lelehan dengan fraksi $F$ dari sumbernya, lalu cairan yang
telah mempertahankan fraksi $f$ dari massanya melalui kristalisasi fraksional, sehingga
untuk unsur dengan $D$ pada residu pelelehan dan $D_c$ pada
mineral yang mengkristal,
\[
C_{\text{lava}} = C_0\,\frac{1}{D + F(1 - D)}\,f^{\,D_c - 1}.
\]
Faktor kedua itu sama untuk setiap unsur yang ditolak oleh mineral
pengkristal, yakni $D_c \approx 0$: faktor itu menskalakan semuanya dengan $1/f$. Rasio
unsur semacam itu karena itu ditentukan di sumbernya, dan lava yang telah kehilangan olivin
tetap mempertahankan La/Yb dari lelehan primernya selagi lantanumnya naik. Membalik $F$
dari suatu konsentrasi juga memerlukan $f$; membalik sumber dari suatu rasio
tidak. Plagioklas adalah pengecualian di antara unsur tanah jarang: ia menyerap europium,
dan pengeluarannya menyisakan anomali europium dari volume Tahun 2.

\section{Membaca pola unsur tanah jarang}\label{sec:b3:trace-element-modelling:patterns}

Lelehan dari lherzolit spinel dan lelehan dari lherzolit garnet berbagi unsur tanah jarang ringan,
yang tidak disimpan oleh kedua residu, dan berbeda pada unsur berat, yang ditahan oleh garnet.
Pada diagram ternormalisasi kondrit, lelehan garnet tampak
curam dan berdasar datar di ujung berat; lelehan spinel lebih landai. Kecuraman
itu meningkat seiring turunnya $F$.

\begin{omfigure}
\begin{tikzpicture}
\begin{axis}[omrecord, width=0.8\linewidth, height=0.5\linewidth, ymode=log,
  xmin=56.5, xmax=71.5, ymin=1, ymax=1000, xtick={57,58,60,62,63,64,65,70,71},
  xticklabels={La,Ce,Nd,Sm,Eu,Gd,Tb,Yb,Lu}, ytick={1,10,100,1000}, yticklabels={1,10,100,1000},
  ylabel={sampel / kondrit CI},
  legend style={font=\scriptsize, at={(0.98,0.98)}, anchor=north east,
  legend cell align=left, legend columns=2}]
% computed in figdata/bachelor-3/melting_models.py
\addplot[omDef, thick, mark=*, mark size=1.3pt] table[x=z, y=sp01] {figdata/out/bachelor-3-melting_models-ree.dat};
\addlegendentry{spinel, $F = 1\,\%$}
\addplot[red!70!black, thick, mark=square*, mark size=1.3pt] table[x=z, y=gt01] {figdata/out/bachelor-3-melting_models-ree.dat};
\addlegendentry{garnet, $F = 1\,\%$}
\addplot[omDef, thick, densely dashed, mark=*, mark size=1.3pt] table[x=z, y=sp05] {figdata/out/bachelor-3-melting_models-ree.dat};
\addlegendentry{spinel, $5\,\%$}
\addplot[red!70!black, thick, densely dashed, mark=square*, mark size=1.3pt] table[x=z, y=gt05] {figdata/out/bachelor-3-melting_models-ree.dat};
\addlegendentry{garnet, $5\,\%$}
\addplot[omDef, thick, densely dotted, mark=*, mark size=1.3pt] table[x=z, y=sp20] {figdata/out/bachelor-3-melting_models-ree.dat};
\addlegendentry{spinel, $20\,\%$}
\addplot[red!70!black, thick, densely dotted, mark=square*, mark size=1.3pt] table[x=z, y=gt20] {figdata/out/bachelor-3-melting_models-ree.dat};
\addlegendentry{garnet, $20\,\%$}
\addplot[black, thin, mark=none] table[x=z, y=pm] {figdata/out/bachelor-3-melting_models-ree.dat};
\addlegendentry{sumber}
\end{axis}
\end{tikzpicture}

\par\smallskip{\small Pola unsur tanah jarang ternormalisasi kondrit dari lelehan tumpak
sumber mantel primitif, dari lherzolit spinel (lingkaran biru) dan
lherzolit garnet (kotak merah), pada tiga fraksi lelehan; sumbernya adalah
garis tipis. Hanya sembilan unsur tanah jarang yang koefisien partisinya tercantum
dalam tabel yang digambarkan.}\label{fig:b3:trace-element-modelling:ree}
\end{omfigure}
% ledger: budahn91:d-ree, ms95:pm-trace, lodders:ci-table, white13:ex73, rfr25:gkr-5gpa

Garnet digantikan oleh spinel sebagai mineral kaya aluminium di mantel pada kedalaman sekitar 60
hingga \qty{90}{km}. Lelehan yang unsur tanah jarang beratnya datar dan rendah, dengan
kemiringan curam dari unsur ringan ke berat, oleh karena itu garnet tertinggal dalam residunya dan
mulai meleleh setidaknya pada kedalaman itu.
% ledger: white13:sp-gt

Suatu pola memuat dua besaran tak diketahui, $F$ dan sumbernya. Dengan memplot dua unsur
satu sama lain akan memisahkan keduanya.

\begin{definition}[Diagram rasio--rasio]\label{def:b3:trace-element-modelling:ratios}
Yang disebut \emph{diagram rasio--rasio}\index{diagram rasio--rasio} memplot, untuk sekumpulan
lava, rasio unsur $H$ yang sangat inkompatibel terhadap unsur $M$ yang kurang
inkompatibel terhadap konsentrasi $H$ (atau terhadap rasio lain). Pelelehan
sebagian, kristalisasi, dan pencampuran menggambar garis dengan bentuk berbeda
di atasnya.
\end{definition}

\begin{proposition}[Garis pelelehan tumpak]\label{prop:b3:trace-element-modelling:line}
Untuk lelehan tumpak dari satu sumber dengan $D_H = 0$ dan $D_M$ konstan, titik-titik
$(C_H, C_H/C_M)$ terletak pada garis lurus
\[
\frac{C_H}{C_M} = \frac{D_M}{C_0^M}\,C_H + \frac{C_0^H(1 - D_M)}{C_0^M}.
\]
\end{proposition}

\begin{proof}
Dengan $D_H = 0$, $C_H = C_0^H/F$, sehingga $1/F = C_H/C_0^H$. Maka
$C_H/C_M = (C_0^H/F)\,(D_M + F(1 - D_M))/C_0^M = C_0^HD_M/(FC_0^M) +
C_0^H(1 - D_M)/C_0^M$, dan suku pertama adalah $(D_M/C_0^M)\,C_H$.
\end{proof}

Kristalisasi fraksional mineral yang menolak kedua unsur menggeser lava
secara horizontal pada diagram ini, karena rasionya tidak berubah; pelelehan menggesernya
sepanjang garis miring yang kemiringannya membesar seiring $D_M$; pencampuran dua lelehan menggambar sebuah
hiperbola. Sumber dengan garnet memiliki $D_{\mathrm{Yb}}$ besar, sehingga garis
pelelehan curam pada La/Yb terhadap La.

\begin{method}[Membalik dua lava]\label{met:b3:trace-element-modelling:invert}
(1) Periksa bahwa lava-lava itu kogenetik dan bahwa $H$ jauh lebih inkompatibel
daripada $M$. (2) Koreksi setiap lava akibat kristalisasi, atau gunakan hanya rasio.
(3) Melalui dua lava atau lebih, cocokkan garis $C_H/C_M = aC_H + b$. (4) Maka
$a = D_M/C_0^M$ dan $b = C_0^H(1 - D_M)/C_0^M$: dengan $C_0^M$ diasumsikan, $D_M =
aC_0^M$ dan $C_0^H = bC_0^M/(1 - D_M)$. (5) Fraksi lelehan tiap lava adalah
$F = C_0^H/C_H$.
\end{method}

\begin{omfigure}
\begin{tikzpicture}
\begin{axis}[omrecord, width=0.78\linewidth, height=0.52\linewidth,
  xmode=log, ymode=log, xmin=2, xmax=60, ymin=2, ymax=100,
  xlabel={La (\unit{ppm})}, ylabel={La/Yb},
  xtick={2,5,10,20,50}, xticklabels={2,5,10,20,50},
  ytick={2,5,10,20,50,100}, yticklabels={2,5,10,20,50,100},
  legend style={font=\scriptsize, at={(0.03,0.97)}, anchor=north west,
  legend cell align=left}]
% computed in figdata/bachelor-3/melting_models.py
\addplot[omDef, thick] table[x=la_sp, y=lay_sp] {figdata/out/bachelor-3-melting_models-traj.dat};
\addlegendentry{lherzolit spinel}
\addplot[red!70!black, thick] table[x=la_gt, y=lay_gt] {figdata/out/bachelor-3-melting_models-traj.dat};
\addlegendentry{lherzolit garnet}
\addplot[omDef, only marks, mark=|, mark size=3pt, forget plot] table[x=la_sp, y=lay_sp] {figdata/out/bachelor-3-melting_models-ticks.dat};
\addplot[red!70!black, only marks, mark=|, mark size=3pt, forget plot] table[x=la_gt, y=lay_gt] {figdata/out/bachelor-3-melting_models-ticks.dat};
\addplot[black, only marks, mark=diamond*, mark size=2.8pt] table[x=la, y=layb] {figdata/out/bachelor-3-melting_models-lavas.dat};
\addlegendentry{lava A dan B}
\node[font=\scriptsize, anchor=north] at (axis cs:9.6,5.2) {A};
\node[font=\scriptsize, anchor=south east] at (axis cs:15.2,36) {B};
\node[font=\scriptsize, omDef, anchor=north west] at (axis cs:36.5,17.5) {1\,\%};
\node[font=\scriptsize, omDef, anchor=north west] at (axis cs:11.3,6.2) {5\,\%};
\node[font=\scriptsize, omDef, anchor=north west] at (axis cs:3.14,2.55) {20\,\%};
\node[font=\scriptsize, red!70!black, anchor=south east] at (axis cs:29.5,66) {1\,\%};
\node[font=\scriptsize, red!70!black, anchor=south east] at (axis cs:10.5,24) {5\,\%};
\node[font=\scriptsize, red!70!black, anchor=south east] at (axis cs:3.09,7.0) {20\,\%};
\end{axis}
\end{tikzpicture}

\par\smallskip{\small Lelehan tumpak dari sumber mantel primitif pada diagram La/Yb
terhadap La, untuk lherzolit spinel dan lherzolit garnet, dengan tanda pada $F = 1$,
2, 5, 10, dan \qty{20}{\percent}, serta dua lava rekaan dari soal akhir
pekan. Setiap lava berada pada satu kurva saja.}\label{fig:b3:trace-element-modelling:ratios}
\end{omfigure}

Pada skala logaritmik garis lurus proposisi itu melengkung, dan
dengan $D_{\mathrm{La}} \ne 0$ garis itu hanya hampir lurus pada skala linear; kurva
pada gambar bersifat eksak. Sumber spinel tidak dapat membentuk lelehan dengan La/Yb di atas
$(C_0^{\mathrm{La}}/D_{\mathrm{La}})/(C_0^{\mathrm{Yb}}/D_{\mathrm{Yb}})$,
sekitar 40 untuk mantel primitif; sumber garnet dapat mencapai sekitar 120.

\section{Latihan}

\begin{exercise}[$\star$]\label{exo:b3:trace-element-modelling:1}
Dengan koefisien mineral samarium, 0.40 dalam klinopiroksen, 0.11 dalam
garnet, 0.016 dalam olivin, ortopiroksen, dan spinel, hitung koefisien
partisi total samarium dalam lherzolit spinel (\qty{12}{\percent}
klinopiroksen) dan dalam lherzolit garnet (\qty{19.5}{\percent} klinopiroksen,
\qty{16.7}{\percent} garnet) dalam bab ini.
\end{exercise}
% ledger: budahn91:d-ree

\begin{exercise}[$\star$]\label{exo:b3:trace-element-modelling:2}
Suatu unsur dengan $D = 0.05$ hadir sebesar \qty{1}{ppm} dalam sumber. Hitung
konsentrasinya dalam lelehan tumpak pada $F = 0.02$ dan $0.10$.
\end{exercise}

\begin{exercise}[$\star$]\label{exo:b3:trace-element-modelling:3}
Berapa pengayaan $C_L/C_0$ terbesar yang dapat diberikan \omterm{thm:b3:trace-element-modelling:batch}{pelelehan tumpak} kepada suatu
unsur dengan $D = 0.01$? Dan kepada suatu unsur dengan $D = 0$ pada
$F = \qty{5}{\percent}$?
\end{exercise}

\begin{exercise}[$\star$]\label{exo:b3:trace-element-modelling:4}
Suatu lava mengandung \qty{30}{ppm} La dan \qty{2.5}{ppm} Yb. Dengan nilai kondrit CI
La 0.232 dan Yb \qty{0.163}{ppm}, hitung La/Yb-nya dan
$(\mathrm{La}/\mathrm{Yb})_N$ ternormalisasi kondrit.
\end{exercise}
% ledger: lodders:ci-table

\begin{exercise}[$\star\star$]\label{exo:b3:trace-element-modelling:5}
Untuk suatu unsur dengan $D = 0.01$, bandingkan residu setelah \qty{5}{\percent}
\omterm{thm:b3:trace-element-modelling:batch}{pelelehan tumpak} dan \omterm{thm:b3:trace-element-modelling:fractional}{pelelehan fraksional}. Apa yang dikandung lelehan fraksional
sesaat pada saat itu?
\end{exercise}

\begin{exercise}[$\star\star$]\label{exo:b3:trace-element-modelling:6}
Dalam contoh kerja dalam bab ini (modus 58, 27, 12, \qty{3}{\percent};
proporsi pelelehan 20, 25, 45, \qty{10}{\percent} untuk olivin,
ortopiroksen, klinopiroksen, dan spinel), lantanum memiliki koefisien
0.0003, 0.002, 0.053, dan 0.01. Hitung $D_0$ dan $P$, lalu pengayaan
lantanum dalam lelehan tumpak pada $F = 0.02$ dan $0.10$, non-modal dan modal.
\end{exercise}
% ledger: white13:ex73

\begin{exercise}[$\star\star$]\label{exo:b3:trace-element-modelling:7}
Pada $F = 0.10$, bandingkan lelehan tumpak dan \omterm{def:b3:trace-element-modelling:fractional}{lelehan fraksional teragregasi} untuk
$D = 0.1$ dan untuk $D = 0.01$. Mana yang lebih berbeda, dan mengapa?
\end{exercise}

\begin{exercise}[$\star\star$]\label{exo:b3:trace-element-modelling:8}
Suatu unsur dengan $D = 0.001$ meleleh secara dinamis dengan $\phi = 0.01$. Hitung
$D'$, pengayaan lelehan pertama yang diekstraksi, dan pengayaan lelehan yang
diekstraksi pada $F = 0.01$; bandingkan dengan \omterm{thm:b3:trace-element-modelling:fractional}{pelelehan fraksional}.
\end{exercise}

\begin{exercise}[$\star\star$]\label{exo:b3:trace-element-modelling:9}
Pada lherzolit spinel bab ini, $D_{\mathrm{La}} = 0.0078$ dan
$D_{\mathrm{Sm}} = 0.062$. Dengan faktor berapa \omterm{thm:b3:trace-element-modelling:batch}{pelelehan tumpak} menaikkan La/Sm di atas
nilai sumber pada $F = 0.01$ dan pada $F = 0.10$? Apa artinya bagi penggunaan
rasio untuk menaksir $F$?
\end{exercise}

\begin{exercise}[$\star\star\star$]\label{exo:b3:trace-element-modelling:10}
Tunjukkan bahwa \omterm{def:b3:trace-element-modelling:fractional}{lelehan fraksional teragregasi} adalah rata-rata dari lelehan sesaat
dari 0 sampai $F$, dan bahwa nilainya menuju $C_0/F$ ketika $(1 - F)^{1/D} \to 0$.
Untuk $D = 0.01$, mulai dari $F$ berapa nilainya berada dalam \qty{1}{\percent} dari $C_0/F$?
\end{exercise}

\begin{exercise}[$\star\star\star$]\label{exo:b3:trace-element-modelling:11}
Dua lava kogenetik P dan Q memiliki La 20 dan \qty{10}{ppm}, serta La/Yb 10 dan 6.
Dengan asumsi lantanum sepenuhnya inkompatibel dan sumber dengan
\qty{0.44}{ppm} Yb, terapkan metode bab ini: tentukan $D_{\mathrm{Yb}}$,
La sumber, dan fraksi lelehan tiap lava.
\end{exercise}

\begin{exercise}[$\star\star\star$]\label{exo:b3:trace-element-modelling:12}
Pada \omterm{def:b3:trace-element-modelling:nonmodal}{pelelehan non-modal} latihan~6, pada fraksi lelehan berapa setiap mineral
habis? Hitung koefisien total residu untuk lantanum pada
$F = 0.2$, dan jelaskan mengapa rumus bab ini tidak dapat dipakai melewati
habisnya pertama.
\end{exercise}

\section{Soal: Dua Lava Pulau Aldren}

\begin{problem}[{Soal akhir pekan --- sumber, kedalaman, dan derajat pelelehan
dua basal rekaan, dari lantanum dan iterbium}]\label{pb:b3:trace-element-modelling:1}
Di Pulau Aldren, pulau vulkanik rekaan, lava A mengandung \qty{9.6}{ppm} La
dan \qty{1.7}{ppm} Yb, lava B \qty{15.6}{ppm} La dan \qty{0.45}{ppm} Yb.
Keduanya diambil sebagai lelehan primer dari sumber mantel primitif dengan
\qty{0.648}{ppm} La dan \qty{0.441}{ppm} Yb (kondrit CI: La 0.232, Yb
\qty{0.163}{ppm}). Koefisien mineral: La 0.058 dalam klinopiroksen dan 0.001
dalam garnet serta dalam olivin, ortopiroksen, dan spinel; Yb 0.58 dalam
klinopiroksen, 4.5 dalam garnet, dan 0.16 dalam olivin, ortopiroksen, dan spinel.
Lherzolit spinel terdiri atas \qty{58}{\percent} olivin, \qty{27}{\percent} ortopiroksen, \qty{12}{\percent}
klinopiroksen, \qty{3}{\percent} spinel; lherzolit garnet terdiri atas \qty{54.1}{\percent} olivin, \qty{9.8}{\percent}
ortopiroksen, \qty{19.5}{\percent} klinopiroksen, \qty{16.7}{\percent} garnet. Pelelehan bersifat modal.
% ledger: ms95:pm-trace, lodders:ci-table, budahn91:d-ree, white13:ex73, rfr25:gkr-5gpa

\medskip\noindent\textbf{Bagian I --- Kedua sumber.}
\begin{enumerate}
\item Hitung $D_{\mathrm{La}}$ dan $D_{\mathrm{Yb}}$ lherzolit spinel.
\item Hitung keduanya untuk lherzolit garnet. Bulatkan keempatnya menjadi dua
angka penting untuk sisa soal.
\item Koefisien mana yang paling berbeda antara kedua sumber, dan mengapa?
Apakah iterbium kompatibel dalam lherzolit garnet?
\item Hitung $(\mathrm{La}/\mathrm{Yb})_N$ sumber ternormalisasi kondrit.
Seperti apa polanya?
\end{enumerate}

\medskip\noindent\textbf{Bagian II --- Pelelehan tumpak lava A.}
\begin{enumerate}[resume]
\item Hitung La dalam lelehan tumpak lherzolit spinel pada $F = 1$, 5, dan
\qty{10}{\percent}.
\item Tentukan fraksi lelehan yang memberikan kandungan lantanum pada lava A.
\item Prediksi kandungan iterbium dan rasio La/Yb-nya, lalu bandingkan dengan lava A.
\item Hitung La dan Yb dalam residu, serta fraksi lantanum dari sumber
yang tersisa di dalamnya.
\item Berapa fraksi lelehan yang kamu peroleh jika lantanum dianggap
sepenuhnya inkompatibel?
\end{enumerate}

\medskip\noindent\textbf{Bagian III --- Pelelehan fraksional lava A.}
\begin{enumerate}[resume]
\item Hitung La dalam lelehan fraksional sesaat pada $F = 0.1$, 1, dan
\qty{6}{\percent}.
\item Hitung La dalam \omterm{def:b3:trace-element-modelling:fractional}{lelehan fraksional teragregasi} pada $F = \qty{6.0}{\percent}$.
\item Tentukan fraksi lelehan saat lelehan teragregasi memiliki kandungan lantanum
lava A. Mengapa hasilnya begitu dekat dengan jawaban pertanyaan~9?
\item Bandingkan kandungan lantanum residu dengan kandungan pada pertanyaan~8.
\item Taksiran $F$ yang mana yang kamu pilih?
\end{enumerate}

\medskip\noindent\textbf{Bagian IV --- Lava B dan garnet.}
\begin{enumerate}[resume]
\item Hitung La/Yb dan $(\mathrm{La}/\mathrm{Yb})_N$ kedua lava.
\item Berapa La/Yb terbesar yang dapat dicapai lelehan tumpak lherzolit
spinel?
\item Pada fraksi lelehan berapa lelehan lherzolit spinel memiliki La/Yb
lava B, dan berapa kandungan lantanum yang dimilikinya? Simpulkan.
\item Tentukan fraksi lelehan lava B dari kandungan lantanumnya untuk
lherzolit garnet, dan prediksi kandungan iterbiumnya.
\item Apa yang dikatakan kedua sumber tentang kedalaman tempat A dan B mulai
meleleh? Mungkinkah lava A berasal dari lherzolit garnet?
\item Hitung kemiringan dan titik potong garis \omterm{thm:b3:trace-element-modelling:batch}{pelelehan tumpak} La/Yb $= a\,
\mathrm{La} + b$ sumber garnet dengan $D_{\mathrm{La}}$ dianggap nol,
dan periksa lava B terhadap garis tersebut.
\item Jika lava A kehilangan \qty{10}{\percent} massanya sebagai olivin sebelum
erupsi, bagaimana La/Yb dan fraksi lelehannya berubah?
\item Jawab pertanyaan pembuka bab ini untuk kedua lava tersebut.
\item Berikan derajat pelelehan lava A.
\end{enumerate}
\end{problem}
